\section{Preference Formation and the Boundary between Risk Choices}
\label{sec:risk_frontier}
The annuity identity identifies which contract changes matter; it does not identify the source of a portfolio reversal. To distinguish an operating-duration incentive from preference formation, fix an initial state and date and partition admissible policies by their \emph{first} risky position. Let $\mathcal P_+$ require a positive first position and $\mathcal P_-$ a nonpositive one. Later actions remain unrestricted within the economic model. In the finite implementation the first position is held over the first rebalancing period, of length $\delta=1/8$, until liquidation. These are nonempty finite policy sets. A continuous-time counterpart must restrict the position over a positive-duration interval, not merely its value at the initial instant; Section \ref{sec:decision_target} makes this distinction explicit. For that counterpart, replace maxima by suprema unless attainment is established. The finite-policy conclusions below concern the declared calendar.

Write
\begin{equation}
 W_\sigma(d,k)=\max_{p\in\mathcal P_\sigma}\{B_p+dA_p-kC_p\},\qquad
 \Delta(d,k)=W_+(d,k)-W_-(d,k),\quad \sigma\in\{+,-\}.
 \label{eq:risk_class_values}
\end{equation}
Thus $\Delta>0$ identifies a strictly positive optimal first position, while $\Delta<0$ rules it out. A zero is an indifference, not necessarily a unique threshold. The economic objects $A_p$ and $C_p$ are discounted duration and adjustment effort, not the first-period duration or the charge $kC_p$.

\begin{proposition}[Duration, adjustment effort, and a risk-switching locus]
\label{prop:risk_frontier}
For finite policy sets, $W_+$ and $W_-$ are convex piecewise-affine functions of $(d,k)$. At a point where the optimizing feature vector in each class is locally constant,
\begin{equation}
 \Delta_d=A_+-A_-,\qquad \Delta_k=-(C_+-C_-).
 \label{eq:risk_derivatives}
\end{equation}
If $\Delta(d^*,k)=0$ and $A_+\ne A_-$ at such a point, the local switching locus has slope
\begin{equation}
 \frac{d d^*(k)}{dk}=\frac{C_+-C_-}{A_+-A_-}.
 \label{eq:risk_frontier_slope}
\end{equation}
More generally, on an interval of $d$ at a fixed $k$, suppose every pair of optimizing policies, one from each class, satisfies $A_+-A_-\geq\kappa>0$. Then $\Delta$ is strictly increasing with slope at least $\kappa$ in the integral sense. Opposite endpoint signs imply exactly one crossing. If a computed difference has uniform error at most $e_++e_-$, any zero of that computed difference lies within $(e_++e_-)/\kappa$ of the true crossing.
\end{proposition}
\begin{proof}
A maximum of finitely many affine functions is convex, Lipschitz on compact sets, and piecewise affine. On a cell with fixed optimizing feature vectors, differentiation gives \eqref{eq:risk_derivatives}; implicit differentiation gives \eqref{eq:risk_frontier_slope}. Across cells the value functions are absolutely continuous. Their derivatives almost everywhere are the features of active optimizers, a finite-policy instance of the envelope argument in \citet{milgrom2002}. Integrating the assumed derivative separation gives
$\Delta(d_2,k)-\Delta(d_1,k)\geq\kappa(d_2-d_1)$. Continuity gives existence and the strict inequality gives uniqueness. Finally, at an approximate zero the absolute true difference is at most $e_++e_-$; the same slope inequality bounds its distance from the true zero.
\end{proof}

The slope ratio gives a contract-equivalent price of differential preference adjustment. If positive-risk policies both keep the agent in operation longer and use more adjustment effort, a higher adjustment cost requires a larger operating benefit before the positive-risk class becomes preferable. If the effort ordering is reversed, the comparative static reverses. The aggregate theorem that optimal effort decreases with $k$ does not determine this ratio: it compares optimal choices across costs, whereas \eqref{eq:risk_frontier_slope} compares two constrained risk classes at the same cost. The duration-separation hypothesis must be established over the claimed interval; a bisection calculation or a derivative at one point does not establish it.

\subsection{A Matched Preference-Formation Counterfactual}
Let $W_\sigma^{\rm adj}$ retain the full adjustment set and let $W_\sigma^0$ restrict $\theta=0$ at every state and date. The state $u$ remains stochastic, with the original volatility, covariance, operating region, and settlement rule. Thus the comparison removes deliberate preference formation, not preference risk, and changes no reward level or financial technology. Define its option value in each risk class by
\begin{equation}
 \Omega_\sigma(d,k)=W_\sigma^{\rm adj}(d,k)-W_\sigma^0(d),\qquad
 \Omega_\sigma\geq0.
 \label{eq:preference_option}
\end{equation}
The zero-adjustment model does not depend on $k$ because its realized adjustment effort is zero.

\begin{proposition}[Risk-specific option values]
\label{prop:option_decomposition}
At every contract,
\begin{equation}
 \Delta^{\rm adj}(d,k)-\Delta^0(d)=\Omega_+(d,k)-\Omega_-(d,k).
 \label{eq:option_decomposition}
\end{equation}
In particular, preference formation reverses a nonpositive-risk ranking, $\Delta^0<0$, precisely when
$\Omega_+-\Omega_->-\Delta^0$. It reverses a positive-risk ranking, $\Delta^0>0$, precisely when
$\Omega_--\Omega_+>\Delta^0$. If $\Delta^0$ has the same strict sign at two contracts but $\Delta^{\rm adj}$ has opposite signs, the change in relative option values identifies the incremental force behind that reversal, rather than the generic annuity identity alone.
\end{proposition}
\begin{proof}
Set inclusion gives nonnegativity in \eqref{eq:preference_option}. Subtract the two definitions of $\Omega_\sigma$ to obtain \eqref{eq:option_decomposition}. The two strict inequalities are then necessary and sufficient for opposite signs. No interior first-order condition or portfolio-distance argument is required.
\end{proof}

The propositions suggest a focused experiment. We independently optimize both risk classes with and without deliberate adjustment, record $(B,A,C)$ for the attaining policies, and compare their value differences. We also vary $k$ while keeping the operating region, covariance, and settlement curvature fixed. The resulting brackets locate an indifference by continuity and opposite signs, but are not reported as a globally unique switching threshold. The full first-action sign certificates from Section \ref{sec:contracts} retain their separate, stronger exclusion interpretation where applicable. Utility units are the same fixed cardinal units as in the original stopped economy; the ratio \eqref{eq:risk_frontier_slope} translates adjustment-cost changes into operating-benefit changes without pretending to estimate either from data.
